% DETECTIVE — Resultados. Documento único que IMPRIME tudo:
%   - Tópicos em PORTUGUÊS (v1, membro -> líder), com métricas de segurança.
%   - Em seguida, os tópicos em INGLÊS (v1 + v2) que você enviou.
% Figuras em Images/: fig_qi_ml, fig_sr_ml, fig_learn_ml, fig_recovery_ml,
%                     fig_qi_ap, fig_sr_ap, fig_attacker_learning, fig_recovery_summary
\documentclass[conference]{IEEEtran}
\usepackage[utf8]{inputenc}
\usepackage[T1]{fontenc}
\usepackage[brazil]{babel}
\usepackage{graphicx}
\usepackage{amsmath}
\usepackage{amssymb}
\usepackage{multirow}
% \ca mantém o conteúdo do comentário no texto (transparente, sem tratamento especial).
\newcommand{\ca}[1]{#1}

\begin{document}

\section{Avaliação de Desempenho}

Avaliamos o DETECTIVE no NS-3 (v3.36) em um cenário de manufatura inteligente
(IIoT). A rede possui $N = 200$ e $250$ nós dispostos aleatoriamente em uma
área de $200 \times 200$~m. Cada execução dura $900$~s, com $150$~s de aquecimento
para a formação dos agrupamentos (\emph{clusters}), sobre IPv6/UDP em uma rede ad
hoc IEEE~802.15.4. Na formação, os nós se agrupam por similaridade de capacidades e
elegem um líder por cluster; cada líder se registra no ponto de acesso (AP) e envia
mensagens periódicas de \emph{heartbeat}.

O modelo de ameaça considera um ataque UDP \emph{flooding} interno (\emph{insider}).
Cinco nós legítimos são comprometidos (cerca de $2\%$ da rede) e passam a inundar o
líder do próprio cluster com tráfego UDP. O ataque começa em $t = 300$~s; cada
atacante envia $50$ pacotes por segundo de $64$~bytes. Um nó que acumula mais de
$500$ pacotes de \emph{flood} satura e sai da rede.

A defesa opera no AP, em ciclos de decisão de $5$~s. A cada ciclo, cada líder
informa ao AP quantos pacotes recebeu de cada origem, e o AP calcula, dentro de cada
cluster, o \emph{Z-score} dessas contagens. A detecção é dupla: (i) um Z-score
relativo, que aponta origens muito acima da média do próprio cluster; e (ii) um
limiar absoluto, que captura o atacante mesmo quando vários atacantes no mesmo
cluster elevam a média e fazem o Z-score colapsar. Para cada origem suspeita, um
agente \emph{Dual-System} com uma Q-table $\epsilon$-greedy escolhe entre duas
ações, \textsc{DoNothing} e \textsc{Quarantine}. O estado é o nível de suspeita,
definido pelo Z-score ($z<2$ limpo, $2 \le z < 4$ baixo, $z \ge 4$ alto). A
recompensa é observacional, sem oráculo do simulador: uma quarentena que deixa o
cluster sem suspeitos é premiada, e uma omissão diante de uma suspeita que se agrava
é penalizada. Antes de isolar um nó, o AP verifica se a sua remoção retiraria alguma
capacidade única do cluster; nesse caso, importa um doador compatível de outro
cluster. Confirmada a decisão, o AP envia uma ordem de quarentena ao líder, que
passa a descartar o tráfego do atacante.

Quando a quarentena é rápida o bastante, o líder sobrevive. Quando o atacante satura
o líder antes da detecção, o líder cai e seus membros ficam órfãos; um pipeline de
recuperação reintegra cada órfão, realocando-o para um cluster vizinho compatível
(\textsc{Reallocate}) ou formando um novo cluster (\textsc{Recluster}). O agente usa
$\alpha_Q = 0.2$, $\gamma_Q = 0.9$ e uma política $\epsilon$-greedy que parte de
$0.15$ e decai por $0.97$ a cada passo até o piso de $0.02$.

Usamos duas métricas de conectividade: o Índice de Qualidade ($QI_t$) e a
Alcançabilidade de Serviço ($SR_t$), medidas antes do ataque, no pior ponto e após
a recuperação. Para avaliar a segurança, comparamos os atacantes plantados
(\emph{ground-truth}) com as decisões do AP e derivamos, da matriz de confusão, a
Acurácia (ACC), a Taxa de Detecção ou \emph{Recall} (DR), a Precisão (PRE), a Taxa
de Falsos Positivos (FPR) e o F1. Também reportamos três tempos: o tempo de detecção
do ataque (ADT), o tempo de mitigação (AMT) e o tempo de recuperação ($T_{rec}$).
Todos os resultados são a média de $35$ execuções por tamanho de rede, com
intervalos de confiança de $95\%$.

\section{Resultados}

A Tabela~\ref{tab:rede} resume a resposta da rede ao ataque. A etapa de formação
produz cerca de $20$ e $23$ líderes, acompanhando a densidade da rede. Sob
ataque, menos de dois líderes por execução caem em média ($1.7$ a $1.9$): a defesa
contém o \emph{flood} cedo, então a maioria dos líderes atacados permanece viva.
Quando um líder cai, ele deixa de $26$ a $34$ órfãos no pior instante, e todos são
recuperados. A política usa apenas \textsc{Reallocate} (cerca de $33$ e $49$
realocações por execução) e nunca \textsc{Recluster}, porque quase sempre há um
cluster vizinho compatível disponível.

\begin{table*}[t]
\centering
\caption{Network response to the insider UDP flooding attack (member $\to$ leader)}
\label{tab:rede}
\setlength{\tabcolsep}{9pt}
\renewcommand{\arraystretch}{1.2}
\begin{tabular}{ccccccc}
\hline
$N$ & Leaders & Sat. leaders & Orphans (peak) & Realloc./run & $SR$ pre (\%) & $SR$ min (\%) \\
\hline
200 & $20.2\pm1.4$ & $1.7\pm0.3$ & $25.9\pm5.5$ & $32.8\pm7.8$ & $98.4\pm0.4$ & $84.0\pm3.1$ \\
250 & $23.3\pm1.0$ & $1.9\pm0.3$ & $34.0\pm5.6$ & $48.8\pm9.8$ & $98.3\pm0.5$ & $83.5\pm2.3$ \\
\hline
\end{tabular}
\end{table*}

A Tabela~\ref{tab:seguranca} traz as métricas de segurança. A detecção é muito
precisa: a Precisão é de $100\%$ e a FPR de $0\%$ em todas as escalas, ou seja, o
DETECTIVE nunca quarentenou um nó legítimo. A Taxa de Detecção fica em torno de
$75\%$ a $79\%$. Ela não é de $100\%$ por dois motivos. Primeiro, alguns atacantes
inundam um líder que o AP nunca registrou, isto é, um nó que age como líder
localmente mas é desconhecido do AP; o AP não vê esse tráfego e não consegue
marcá-lo. Segundo, poucos líderes caem já na primeira janela de detecção, antes que
o Z-score consiga separar o atacante dos membros normais. A alta Acurácia (acima de
$99\%$) reflete a base majoritária de nós normais corretamente não isolados.

\begin{table*}[t]
\centering
\caption{DETECTIVE security metrics (member $\to$ leader)}
\label{tab:seguranca}
\setlength{\tabcolsep}{9pt}
\renewcommand{\arraystretch}{1.2}
\begin{tabular}{cccccccc}
\hline
$N$ & ACC (\%) & DR / Recall (\%) & PRE (\%) & FPR (\%) & F1 (\%) & ADT (s) & $T_{rec}$ (s) \\
\hline
200 & $99.49$ & $79.43$ & $100.0$ & $0.00$ & $88.54$ & $5.9$ & $5.4$ \\
250 & $99.49$ & $74.55$ & $100.0$ & $0.00$ & $85.42$ & $5.2$ & $5.6$ \\
\hline
\end{tabular}
\end{table*}

A Tabela~\ref{tab:confusao} mostra a matriz de confusão agregada (por nó, somando as
duas escalas): a defesa isola os atacantes (TP) \emph{sem gerar nenhum falso
positivo} (FP $=0$), e os erros são apenas atacantes não detectados (FN), coerente
com a Precisão de $100\%$ e a FPR de $0\%$.

\begin{table}[t]
\centering
\caption{Confusion matrix (aggregate, member $\to$ leader)}
\label{tab:confusao}
\setlength{\tabcolsep}{8pt}
\renewcommand{\arraystretch}{1.3}
\begin{tabular}{cc|cc}
\multicolumn{2}{c|}{} & \multicolumn{2}{c}{Predicted} \\
\multicolumn{2}{c|}{} & Attacker & Normal \\
\hline
\multirow{2}{*}{Actual} & Attacker & $262$ (TP) & $78$ (FN) \\
                        & Normal   & $0$ (FP)   & $14910$ (TN) \\
\hline
\end{tabular}
\end{table}

Quanto aos tempos, o ataque é detectado em cerca de $5$ a $6$~s após seu início, ou
seja, no primeiro ciclo de decisão do AP após $t=300$~s. A mitigação segue quase no
mesmo instante da detecção (tempo de mitigação abaixo de $1$~s), pois a decisão de
quarentena é tomada no mesmo ciclo. A rede recupera a alcançabilidade em torno de
$5$ a $6$~s, retornando ao patamar anterior ao ataque em um a dois ciclos. Somando
ao instante do ataque, a recuperação completa ($t_{\text{rec}}^{\text{abs}} = 300 +
\text{ADT} + T_{rec}$) ocorre por volta de $t \approx 311$~s ($311.3$~s para
$N=200$ e $310.8$~s para $N=250$).

A Figura~\ref{fig:qisr} alinha o $QI_t$ e o $SR_t$ ao instante do ataque ($t=0$).
Ambas as métricas ficam próximas de $98\%$ até cerca de $+15$~s, quando os primeiros
líderes caem. O $SR$ cai para cerca de $84\%$ e o $QI$ para a casa dos $90\%$, e
ambos retornam a $96$--$97\%$ por volta de $+30$~s. Eles não voltam totalmente ao
valor inicial: a diferença de cerca de dois pontos corresponde aos líderes perdidos
por saturação, que não são reconstruídos, mas cujos membros seguem atendidos por
líderes vizinhos.

\begin{figure*}[t]
\centering
\includegraphics[width=0.49\textwidth]{Images/fig_qi_ml.pdf}
\hfill
\includegraphics[width=0.49\textwidth]{Images/fig_sr_ml.pdf}
\caption{Índice de Qualidade e Alcançabilidade de Serviço em torno do ataque
($t=0$). As linhas são a média entre execuções; a faixa é $\pm1$ desvio-padrão.}
\label{fig:qisr}
\end{figure*}

A Figura~\ref{fig:learn} acompanha o agente atacante ao longo das $35$ execuções. O
$Q(\textsc{Quarantine})$ sobe de perto de zero para cerca de $10$--$17$ nas
primeiras execuções e se mantém, ou seja, o agente aprende cedo que quarentenar uma
origem marcada é melhor do que não agir. A taxa de exploração $\epsilon$ cai até o
piso de $0.02$ entre as rodadas $20$ e $25$ em todas as escalas; a partir daí o
agente deixa de explorar e mantém a quarentena.

\begin{figure*}[t]
\centering
\includegraphics[width=0.9\textwidth]{Images/fig_learn_ml.pdf}
\caption{Agente atacante ao longo das execuções cumulativas: (A)
$Q(\textsc{Quarantine})$ e (B) taxa de exploração $\epsilon$.}
\label{fig:learn}
\end{figure*}

A Figura~\ref{fig:recovery} resume o $SR$ antes do ataque, no pior ponto e após a
recuperação, para cada tamanho de rede. Em todas as escalas, o DETECTIVE mantém o
ataque sob controle, isola os atacantes sem falsos positivos e devolve o serviço a
um nível próximo do original. A camada de detecção evita a maior parte do dano, e a
camada de recuperação trata o que resta.

\begin{figure}[t]
\centering
\includegraphics[width=\columnwidth]{Images/fig_recovery_ml.pdf}
\caption{Alcançabilidade de Serviço antes do ataque, no pior ponto e após a
recuperação, para cada tamanho de rede.}
\label{fig:recovery}
\end{figure}

% =====================================================================
% Tópicos em INGLÊS (v1 + v2) — conforme enviado.
% =====================================================================

\section{Performance Evaluation}

We evaluated DETECTIVE in NS-3 (v3.36) in a smart-manufacturing IIoT setting.
The network has $N = 200$ and $250$
nodes placed at random over a $200 \times 200$~m area. Each run lasts $900$~s,
with a $150$~s warm-up for cluster formation, over IPv6/UDP on an IEEE~802.15.4 ad
hoc network. \ca{O modelo de ameaça  considera o ataque UDP flooding.... e atua de duas maneiras.... }

\ca{Os ataques acontecem no instante t = 300, com x porcento de atacantes}

(ser mais direto The stress is an insider UDP flooding attack, and we test it in two forms. In the first, compromised members flood their own cluster leader (member\,$\to$\,leader). In the second, compromised leaders flood the Access Point (AP) directly (leader\,$\to$\,AP). In both cases, five attackers start at $t=300$~s, each one sending $50$ packets per second of $64$~bytes. A node that receives more than $500$ flood packets saturates and goes down.

The AP monitors looks at two kinds of information: the traffic each leader reports and the traffic it receives itself. In the member\,$\to$\,leader attack, each leader tells the AP how many packets it got from each source, and the AP flags a source with a Z-score inside the cluster. In the leader\,$\to$\,AP attack, the AP looks at the flood it receives directly. No normal node floods the AP, so any source above a small threshold is an attacker, and the AP checks this every second so it can act before becoming saturated. \ca{In both cases a small Dual-System agent selects one of two actions, \textsc{DoNothing} and \textsc{Quarantine}. A quarantined member is dropped
by its leader. A quarantined leader is removed from the topology, and its members
become orphans that an orphan-recovery pipeline handles with \textsc{Reallocate}
(move the orphan to a nearby cluster) or \textsc{Recluster} (build a new cluster
from the orphans). The attacker agent uses $\alpha_Q = 0.2$, $\gamma_Q = 0.9$, and
an $\epsilon$-greedy policy that starts at $0.15$ and decays by $0.97$ per step down
to $0.02$.}

We use two connectivity metrics: the Quality Index ($QI_t$) and the
Service Reachability ($SR_t$). We measure them before the attack, at the worst
point, and after recovery. We also report the quarantine recall, which is the share
of attackers the AP isolates, and the learning of the attacker agent through
$Q(\textsc{Quarantine})$ and $\epsilon$. All results are the average of $35$ runs
per network size, with $95\%$ confidence intervals.

\section{Results}

Table~\ref{tab:detective} shows the behavior of DETECTIVE under the two attacks.
The two cases are different in how hard they are to detect and in how much damage
they do. In the member\,$\to$\,leader attack, recall is around $75\%$ to $79\%$.
It is not $100\%$ for two reasons. First, some attackers flood a leader
that the AP never registered, that is, a node that acts as a leader locally but is
unknown to the AP; the AP does not see this traffic, so it cannot flag these
attackers. Second, a few leaders go down in the first detection window, before the
Z-score can tell the attacker apart from the normal members. In the
leader\,$\to$\,AP attack, recall is $100\%$ in every
run. Here the AP sees the flood directly, so no attacker can hide, and the AP is
never lost to saturation.

\begin{table*}[t]
\centering
\caption{DETECTIVE under the two insider flooding attacks}
\label{tab:detective}
\setlength{\tabcolsep}{10pt}
\renewcommand{\arraystretch}{1.2}
\begin{tabular}{cccccccc}
\hline
$N$ & Attack & Leaders & Orphans (peak) & Recall (\%) & $SR$ pre (\%) & $SR$ min (\%) & $SR$ final (\%) \\
\hline
200 & member\,$\to$\,leader & $20.2\pm1.4$ & $25.9\pm5.5$ & $79.4$ & $98.4\pm0.4$ & $84.0\pm3.1$ & $96.8\pm0.6$ \\
200 & leader\,$\to$\,AP      & $19.5\pm0.8$ & $48.2\pm7.3$ & $100$  & $98.6\pm0.3$ & $72.0\pm3.6$ & $96.1\pm0.3$ \\
\hline
250 & member\,$\to$\,leader & $23.3\pm1.0$ & $34.0\pm5.6$ & $74.5$ & $98.3\pm0.5$ & $83.5\pm2.3$ & $96.9\pm0.5$ \\
250 & leader\,$\to$\,AP      & $23.4\pm1.2$ & $49.7\pm8.2$ & $100$  & $98.0\pm0.6$ & $76.2\pm3.3$ & $96.1\pm0.6$ \\
\hline
\end{tabular}
\end{table*}

The two attacks also differ in how the network recovers. In the
member\,$\to$\,leader attack, fewer than two leaders per run go down on average
($1.7$ to $1.9$). The attack is stopped early, so most attacked leaders stay alive.
When a leader does go down, it leaves $26$ to $34$ orphans, and all of them are
moved to other clusters. The agent uses only \textsc{Reallocate} (about $33$ and
$49$ moves per run) and never \textsc{Recluster}, because a nearby working
cluster is always available. In the leader\,$\to$\,AP attack, each attack removes
five leaders at once, which leaves about $48$ to $50$ orphans. All of them are
recovered as well, also almost exclusively through \textsc{Reallocate}.

Figures~\ref{fig:qisr_ml} and \ref{fig:qisr_ap} line up $QI_t$ and $SR_t$ with the
attack time ($t=0$). In the member\,$\to$\,leader attack (Figure~\ref{fig:qisr_ml}),
both metrics stay near $98\%$ until about $+15$~s, when the first leaders go down.
$SR$ falls to about $84\%$ and $QI$ to the low nineties, then goes back to
$96$--$97\%$ by $+30$~s. In the leader\,$\to$\,AP attack
(Figure~\ref{fig:qisr_ap}), the drop is sharper and deeper, because five leaders and
their members leave at the same time: $SR$ falls to $72$--$76\%$ right at $t=0$, and
then returns to about $96\%$ within one or two AP cycles. In both attacks the
metrics do not return all the way to the start. The gap of about two points is
caused by the leaders that were removed, which are not rebuilt. The drop is deeper
for smaller networks, because losing a fixed number of leaders is a larger share of
a small network.

\begin{figure*}[t]
\centering
\includegraphics[width=0.49\textwidth]{Images/fig_qi_ml.pdf}
\hfill
\includegraphics[width=0.49\textwidth]{Images/fig_sr_ml.pdf}
\caption{member\,$\to$\,leader attack: Quality Index and Service Reachability around
the attack ($t=0$). Lines are the average over runs; the band is $\pm1$ standard
deviation.}
\label{fig:qisr_ml}
\end{figure*}

\begin{figure*}[t]
\centering
\includegraphics[width=0.49\textwidth]{Images/fig_qi_ap.pdf}
\hfill
\includegraphics[width=0.49\textwidth]{Images/fig_sr_ap.pdf}
\caption{leader\,$\to$\,AP attack: Quality Index and Service Reachability around the
attack ($t=0$).}
\label{fig:qisr_ap}
\end{figure*}

Figure~\ref{fig:learn_en} follows the attacker agent over the $35$ runs.
$Q(\textsc{Quarantine})$ rises from near zero to about $10$--$17$ in the first runs
and stays there, so the agent learns early that quarantining a flagged node is
better than doing nothing. The exploration rate $\epsilon$ drops to its floor of
$0.02$ between rounds $20$ and $25$ in all sizes. After that the agent stops
exploring and keeps quarantining.

\begin{figure*}[t]
\centering
\includegraphics[width=0.9\textwidth]{Images/fig_attacker_learning.pdf}
\caption{Attacker agent over cumulative runs: (A) $Q(\textsc{Quarantine})$ and
(B) exploration rate $\epsilon$.}
\label{fig:learn_en}
\end{figure*}

Figure~\ref{fig:recovery_en} sums up the $SR$ before the attack, at the worst point,
and after recovery. It shows the same result in one view: a short drop followed by a
recovery to within two points of the start. Across both attacks and all sizes,
DETECTIVE keeps the network under control, isolates the attackers, and brings
service back to near its original level. The detection layer stops most of the
damage, and the recovery layer handles what is left.

\begin{figure}[t]
\centering
\includegraphics[width=\columnwidth]{Images/fig_recovery_summary.pdf}
\caption{Service Reachability before the attack, at the worst point, and after
recovery, for each network size.}
\label{fig:recovery_en}
\end{figure}

\end{document}
